% GENERATED FILE. Edit canonical lessons, then run npm run generate:books.
% canonical-source-hash: fnv1a64:97edee03e0520918
% canonical-lessons: GU-C324-samrajya, GU-C324-samrat, GU-C324-vans, GU-C324-kabjo, GU-C324-uthlavvu, GU-R324-review-317-324

\chapter{An Empire Falls}
\label{ch:gu-an-empire-falls}
\hlchaptermodality{324}

\begin{chapteropening}
\textbf{By the end of this chapter:} \emph{I can tell how an empire grew and how its last ruler fell.}

Complete the last lesson of chapter 324: I can tell how an empire grew and how its last ruler fell.
\end{chapteropening}

\section[\texorpdfstring{\gu{સામ્રાજ્ય} (sāmrājya)}{સામ્રાજ્ય (sāmrājya)}]{\gu{સામ્રાજ્ય} (sāmrājya) — an empire}

\label{lesson:GU-C324-samrajya}



\begin{quote}
Before the new one: say the Gujarati for enmity, then the Gujarati for to challenge.
\end{quote}

\subsection*{You'll want to know: \gu{સામ્રાજ્ય}}
\textbf{\gu{સામ્રાજ્ય}} — \emph{sāmrājya} — \textquotedblleft{}an empire\textquotedblright{}.

\begin{quote}
\gu{આ} \gu{સામ્રાજ્ય} \gu{દરિયાથી} \gu{પર્વત} \gu{સુધી} \gu{હતું}.
\end{quote}

\begin{quote}
\emph{ā sāmrājya dariyāthī parvat sudhī hatũ.}
\end{quote}

\begin{quote}
This empire stretched from the sea to the mountains.
\end{quote}

A neuter noun.

\begin{cousinweb}
\textbf{\gu{સામ્રાજ્ય}} (\emph{sāmrājya}) is a tatsama, Sanskrit \emph{sāmrājya}, \textquotedblleft{}complete sovereignty, an empire\textquotedblright{} (Monier-Williams), built on \emph{samrāj}, a supreme ruler.
\end{cousinweb}

\begin{grammarlens}[title={What you've built}]
One more word for telling what happened.
\end{grammarlens}

\subsection*{Your turn}
\begin{itemize}
\raggedright
  \item \emph{Say it:} \emph{sāmrājya}
  \item \emph{Say it:} \emph{sāmrājya}, once more
  \item \emph{Say it:} \emph{ā sāmrājya dariyāthī parvat sudhī hatũ.} — This empire stretched from the sea to the mountains.
  \item \emph{Recall:} say the Gujarati for enmity, then the Gujarati for to challenge, then say \emph{sāmrājya} again
\end{itemize}

\begin{tcolorbox}[breakable,colback=teal!4,colframe=teal!35!black,title={Before you move on}]
What does \gu{સામ્રાજ્ય} mean? (\textquotedblleft{}An empire\textquotedblright{}.) Say it once more.
\end{tcolorbox}
\section[\texorpdfstring{\gu{સમ્રાટ} (samrāṭ)}{સમ્રાટ (samrāṭ)}]{\gu{સમ્રાટ} (samrāṭ) — an emperor}

\label{lesson:GU-C324-samrat}



\begin{quote}
Before the new one: say the Gujarati for to challenge, then the Gujarati for an empire.
\end{quote}

\subsection*{You'll want to know: \gu{સમ્રાટ}}
\textbf{\gu{સમ્રાટ}} — \emph{samrāṭ} — \textquotedblleft{}an emperor\textquotedblright{}.

\begin{quote}
\gu{સમ્રાટ} \gu{પાસે} \gu{મોટી} \gu{સેના} \gu{હતી}.
\end{quote}

\begin{quote}
\emph{samrāṭ pāse moṭī senā hatī.}
\end{quote}

\begin{quote}
The emperor had a large army.
\end{quote}

A masculine noun.

\begin{cousinweb}
\textbf{\gu{સમ્રાટ}} (\emph{samrāṭ}) is a tatsama, Sanskrit \emph{samrāj}, \textquotedblleft{}a supreme ruler\textquotedblright{} (Monier-Williams), whose nominative is \emph{samrāṭ}: a king above kings. His realm is the \textbf{\gu{સામ્રાજ્ય}}.
\end{cousinweb}

\begin{grammarlens}[title={What you've built}]
One more word for telling what happened.
\end{grammarlens}

\subsection*{Your turn}
\begin{itemize}
\raggedright
  \item \emph{Say it:} \emph{samrāṭ}
  \item \emph{Say it:} \emph{samrāṭ}, once more
  \item \emph{Say it:} \emph{samrāṭ pāse moṭī senā hatī.} — The emperor had a large army.
  \item \emph{Recall:} say the Gujarati for to challenge, then the Gujarati for an empire, then say \emph{samrāṭ} again
\end{itemize}

\begin{tcolorbox}[breakable,colback=teal!4,colframe=teal!35!black,title={Before you move on}]
What does \gu{સમ્રાટ} mean? (\textquotedblleft{}An emperor\textquotedblright{}.) Say it once more.
\end{tcolorbox}
\section[\texorpdfstring{\gu{વંશ} (vanś)}{વંશ (vanś)}]{\gu{વંશ} (vanś) — a dynasty; a lineage}

\label{lesson:GU-C324-vans}



\begin{quote}
Before the new one: say the Gujarati for an empire, then the Gujarati for an emperor.
\end{quote}

\subsection*{You'll want to know: \gu{વંશ}}
\textbf{\gu{વંશ}} — \emph{vanś} — \textquotedblleft{}a dynasty; a lineage\textquotedblright{}.

\begin{quote}
\gu{આ} \gu{વંશનો} \gu{છેલ્લો} \gu{રાજા} \gu{ખૂબ} \gu{નબળો} \gu{હતો}.
\end{quote}

\begin{quote}
\emph{ā vanśno chello rājā khūb nabḷo hato.}
\end{quote}

\begin{quote}
The last king of this dynasty was very weak.
\end{quote}

A masculine noun.

\begin{cousinweb}
\textbf{\gu{વંશ}} (\emph{vanś}) is a tatsama, Sanskrit \emph{vaṁśa}, \textquotedblleft{}a bamboo cane\textquotedblright{}, and then \textquotedblleft{}a family line\textquotedblright{}, one joint after another (CDIAL 11175). It is masculine: \textbf{\gu{આ} \gu{વંશનો} \gu{રાજા}}.
\end{cousinweb}

\begin{grammarlens}[title={What you've built}]
One more word for telling what happened.
\end{grammarlens}

\subsection*{Your turn}
\begin{itemize}
\raggedright
  \item \emph{Say it:} \emph{vanś}
  \item \emph{Say it:} \emph{vanś}, once more
  \item \emph{Say it:} \emph{ā vanśno chello rājā khūb nabḷo hato.} — The last king of this dynasty was very weak.
  \item \emph{Recall:} say the Gujarati for an empire, then the Gujarati for an emperor, then say \emph{vanś} again
\end{itemize}

\begin{tcolorbox}[breakable,colback=teal!4,colframe=teal!35!black,title={Before you move on}]
What does \gu{વંશ} mean? (\textquotedblleft{}A dynasty; a lineage\textquotedblright{}.) Say it once more.
\end{tcolorbox}
\section[\texorpdfstring{\gu{કબજો} (kabjo)}{કબજો (kabjo)}]{\gu{કબજો} (kabjo) — control; possession}

\label{lesson:GU-C324-kabjo}



\begin{quote}
Before the new one: say the Gujarati for an emperor, then the Gujarati for a dynasty.
\end{quote}

\subsection*{You'll want to know: \gu{કબજો}}
\textbf{\gu{કબજો}} — \emph{kabjo} — \textquotedblleft{}control; possession\textquotedblright{}.

\begin{quote}
\gu{દુશ્મને} \gu{કિલ્લા} \gu{પર} \gu{કબજો} \gu{કર્યો}.
\end{quote}

\begin{quote}
\emph{duśmane killā par kabjo karyo.}
\end{quote}

\begin{quote}
The enemy took control of the fort.
\end{quote}

A masculine noun.

\begin{cousinweb}
\textbf{\gu{કબજો}} (\emph{kabjo}) is usually traced to Persian \emph{qabza}, \textquotedblleft{}a grasp, a hold\textquotedblright{}, a word Persian took from Arabic. A place is taken with \textbf{\gu{પર}}: \textbf{\gu{કિલ્લા} \gu{પર} \gu{કબજો} \gu{કર્યો}}.
\end{cousinweb}

\begin{grammarlens}[title={What you've built}]
One more word for telling what happened.
\end{grammarlens}

\subsection*{Your turn}
\begin{itemize}
\raggedright
  \item \emph{Say it:} \emph{kabjo}
  \item \emph{Say it:} \emph{kabjo}, once more
  \item \emph{Say it:} \emph{duśmane killā par kabjo karyo.} — The enemy took control of the fort.
  \item \emph{Recall:} say the Gujarati for an emperor, then the Gujarati for a dynasty, then say \emph{kabjo} again
\end{itemize}

\begin{tcolorbox}[breakable,colback=teal!4,colframe=teal!35!black,title={Before you move on}]
What does \gu{કબજો} mean? (\textquotedblleft{}Control; possession\textquotedblright{}.) Say it once more.
\end{tcolorbox}
\section[\texorpdfstring{\gu{ઉથલાવવું} (uthlāvvũ)}{ઉથલાવવું (uthlāvvũ)}]{\gu{ઉથલાવવું} (uthlāvvũ) — to overthrow}

\label{lesson:GU-C324-uthlavvu}



\begin{quote}
Before the new one: say the Gujarati for a dynasty, then the Gujarati for control.
\end{quote}

\subsection*{You'll want to know: \gu{ઉથલાવવું}}
\textbf{\gu{ઉથલાવવું}} — \emph{uthlāvvũ} — \textquotedblleft{}to overthrow\textquotedblright{}.

\begin{quote}
\gu{લોકોએ} \gu{નબળા} \gu{રાજાને} \gu{ઉથલાવ્યો}.
\end{quote}

\begin{quote}
\emph{lokoe nabḷā rājāne uthlāvyo.}
\end{quote}

\begin{quote}
The people overthrew the weak king.
\end{quote}

Past, with \gu{તેણે} (he or she …): \gu{તેણે} \gu{રાજાને} \gu{ઉથલાવ્યો}, \emph{teṇe rājāne uthlāvyo}.

\begin{cousinweb}
\textbf{\gu{ઉથલાવવું}} (\emph{uthlāvvũ}) is the causative of \emph{uthalvũ}, to be upset, which CDIAL lists under \emph{*utthalati}, a form it reconstructs (CDIAL 1805): to turn something over, and so to bring down a ruler.
\end{cousinweb}

\begin{grammarlens}[title={What you've built}]
One more word for telling what happened.
\end{grammarlens}

\subsection*{Your turn}
\begin{itemize}
\raggedright
  \item \emph{Say it:} \emph{uthlāvvũ}
  \item \emph{Say it:} \emph{uthlāvvũ}, once more
  \item \emph{Say it:} \emph{lokoe nabḷā rājāne uthlāvyo.} — The people overthrew the weak king.
  \item \emph{Recall:} say the Gujarati for a dynasty, then the Gujarati for control, then say \emph{uthlāvvũ} again
\end{itemize}

\begin{tcolorbox}[breakable,colback=teal!4,colframe=teal!35!black,title={Before you move on}]
What does \gu{ઉથલાવવું} mean? (\textquotedblleft{}To overthrow\textquotedblright{}.) Say it once more.
\end{tcolorbox}
\section[(dialogue)]{Looking back: from Pages of History to An Empire Falls}

\label{lesson:GU-R324-review-317-324}



\begin{quote}
Say the Gujarati for history and a century.
\end{quote}

\subsection*{Your turn}
Say these aloud:

\begin{itemize}
\raggedright
  \item history, a century, ancient, the past, to be born
  \item a king, a queen, a palace, a royal court, a throne
  \item a fort, ruins, a memorial, a statue, to settle
  \item a war, an army, a soldier, an enemy, an invasion
  \item a sword, a shield, a weapon, a cannon, to pierce
  \item an attack, a battle, a warrior, to slay, to fight on
  \item a martyr, a sacrifice, a victory, enmity, to challenge
  \item an empire, an emperor, a dynasty, control, to overthrow
\end{itemize}

\begin{tcolorbox}[breakable,colback=teal!4,colframe=teal!35!black,title={Before you move on}]
History? (\textbf{\gu{ઇતિહાસ}}, \emph{itihās}.) A century? (\textbf{\gu{સદી}}, \emph{sadī}.) Ancient? (\textbf{\gu{પ્રાચીન}}, \emph{prācīn}.) The past? (\textbf{\gu{ભૂતકાળ}}, \emph{bhūtkāḷ}.) To be born? (\textbf{\gu{જન્મવું}}, \emph{janmavũ}.) A king? (\textbf{\gu{રાજા}}, \emph{rājā}.) A queen? (\textbf{\gu{રાણી}}, \emph{rāṇī}.) A palace? (\textbf{\gu{મહેલ}}, \emph{mahel}.) A royal court? (\textbf{\gu{દરબાર}}, \emph{darbār}.) A throne? (\textbf{\gu{ગાદી}}, \emph{gādī}.) A fort? (\textbf{\gu{કિલ્લો}}, \emph{killo}.) Ruins? (\textbf{\gu{ખંડેર}}, \emph{khanḍer}.) A memorial? (\textbf{\gu{સ્મારક}}, \emph{smārak}.) A statue? (\textbf{\gu{મૂર્તિ}}, \emph{mūrti}.) To settle? (\textbf{\gu{વસવું}}, \emph{vasvũ}.) A war? (\textbf{\gu{યુદ્ધ}}, \emph{yuddh}.) An army? (\textbf{\gu{સેના}}, \emph{senā}.) A soldier? (\textbf{\gu{સૈનિક}}, \emph{sainik}.) An enemy? (\textbf{\gu{દુશ્મન}}, \emph{duśman}.) An invasion? (\textbf{\gu{આક્રમણ}}, \emph{ākramaṇ}.) A sword? (\textbf{\gu{તલવાર}}, \emph{talvār}.) A shield? (\textbf{\gu{ઢાલ}}, \emph{ḍhāl}.) A weapon? (\textbf{\gu{હથિયાર}}, \emph{hathiyār}.) A cannon? (\textbf{\gu{તોપ}}, \emph{top}.) To pierce? (\textbf{\gu{વીંધવું}}, \emph{vī̃dhvũ}.) An attack? (\textbf{\gu{હુમલો}}, \emph{humlo}.) A battle? (\textbf{\gu{લડાઈ}}, \emph{laḍāī}.) A warrior? (\textbf{\gu{યોદ્ધો}}, \emph{yoddho}.) To slay? (\textbf{\gu{હણવું}}, \emph{haṇvũ}.) To fight on? (\textbf{\gu{ઝઝૂમવું}}, \emph{jhajhūmvũ}.) A martyr? (\textbf{\gu{શહીદ}}, \emph{śahīd}.) A sacrifice? (\textbf{\gu{બલિદાન}}, \emph{balidān}.) A victory? (\textbf{\gu{વિજય}}, \emph{vijay}.) Enmity? (\textbf{\gu{વેર}}, \emph{ver}.) To challenge? (\textbf{\gu{પડકારવું}}, \emph{paḍkārvũ}.) An empire? (\textbf{\gu{સામ્રાજ્ય}}, \emph{sāmrājya}.) An emperor? (\textbf{\gu{સમ્રાટ}}, \emph{samrāṭ}.) A dynasty? (\textbf{\gu{વંશ}}, \emph{vanś}.) Control? (\textbf{\gu{કબજો}}, \emph{kabjo}.) To overthrow? (\textbf{\gu{ઉથલાવવું}}, \emph{uthlāvvũ}.)
\end{tcolorbox}
